视觉的根本目的并不是输出一串标签，而是从不完整、受噪声污染的输入中恢复一个足够稳定、可操作、可用于定位、预测与想象的外部世界表征。受心理物理学、系统神经科学与 psychometrics 中“通过密集行为探测反推内部表征结构”的思路启发，本文提出 \textbf{Ordinary-Bench}：一个通过受控 3D 场景、序数空间关系问答与场景级重建来评估视觉语言模型（VLM）空间信念的 benchmark。本文的核心观察是：既然充分且准确的偏序关系可以近似重建场景，那么对 VLM 进行大规模 QRR/TRR 探测，本质上就是在读取它“脑海中的场景”是否存在、是否一致、是否可复原。基于当前 540 个场景、200000+ QA 对的实验，我们发现模型在空间信念层面的能力远弱于表面问答分数所暗示的水平：单图条件下 QRR 准确率约为 70\%，多视角输入反而下降到接近随机猜测，TRR 长期停留在随机水平，而基于 200K+ QA 对的微调或强化学习也没有带来显著改善。与生成式 world model 或依赖工具使用、长思维链和 test-time search 的 visual CoT 方法不同，Ordinary-Bench 评估的不是模型能否借助额外流程把题做对，而是其回答本身能否定义一个对象级、代码级、可矢量化重建的场景结构。由此，我们将评测问题从“模型答对了多少空间题”推进到一个更严格的问题：\emph{模型是否真正形成了一个可被重建、可被想象、可被对齐到真实世界的空间场景信念}，并为后续的场景矢量化重建、定位与结构化视觉记忆提供可操作的评估路径。
